\documentclass[11pt,a4paper]{article}
\usepackage[utf8]{inputenc}
\usepackage[ngerman]{babel}
\usepackage{amsmath,amssymb,amsthm}
\usepackage{geometry}
\geometry{margin=2.5cm}

\title{Geometrische Fundierung und Topologische Quantisierung der Gravitationsschnittstelle im EABC-Modell}
\author{Thomas Hoffbauer}
\date{\today}

\begin{document}

\maketitle

\section{Epistemologische Baseline und methodische Demarkation}

Ein direktes, unkonditioniertes Herleiten der klassischen vierdimensionalen Einsteinschen Feldgleichungen (wie der Schwarzschild- oder Kerr-Metrik) aus dem euklidischen, 8-dimensionalen Faktorwortraum ist im aktuellen Modellzustand nicht ausformuliert. Das Framework deklariert dies explizit als geometrisch-physikalische Grenzüberschreitung und vermeidet damit systematische Kategorienfehler. Der oktonionische Zustandsraum $\mathbb{O} \cong \mathbb{R}^8$ spannt von Natur aus einen euklidischen, Riemannschen Raum mit einer strikt positiv definiten Metrik auf. Ihn eins zu eins mit der physikalischen, lorentzschen Raumzeit der Signatur $(3,1)$ gleichzusetzen, ist unzulässig. 

Der Übergang erfordert zwingend eine kontrollierte geometrische Signaturspaltung und Dimensionsreduktion (etwa über die Hopf-Faserung $S^7 \to S^4$ oder die Zorn-Vektor-Spaltung), welche im Framework methodisch sauber als offene \emph{Bridge Conjecture} ausgewiesen ist, um Zirkelschlüsse strikt zu umgehen.

\section{Die Vakuum-Feldgleichung in acht Dimensionen: $\operatorname{Ric}(g) = 0$ und Holonomie}

Auf der differentialgeometrischen Ebene des Levi-Civita-Suchprogramms führt das System zu einer harten mathematischen Fessel. Es wird gefordert, dass die aus der Oktonionen-Multiplikation stammende Cayley-Vierform $\Phi$ bezüglich der Levi-Civita-Verbindung parallel ist ($\nabla^g \Phi = 0$). Nach dem mathematischen Satz von Berger reduziert dies die Holonomie der 8-dimensionalen Mannigfaltigkeit auf $\operatorname{Spin}(7)$, was wegen der Torsionsfreiheit des Programms die strikte Ricci-Flachheit ($\operatorname{Ric}(g) = 0$) erzwingt.

Dies fungiert als geometrischer Selektionsfilter für die physikalisch zulässigen Metriken des Hilbertraums. Metrik-Konfigurationen, die diese Gleichung verletzen -- namentlich die unphysikalischen 21-dimensionalen Weyl-Rauschmoden --, erfahren an den Gitter-Potentialbarrieren eine unerbittliche Dämpfung und werden energetisch komplett aus dem physikalischen Raum verbannt.

\section{Die Holografische Kopplungsbedingung an der Planck-Skala}

An der Schnittstelle zur Thermodynamik schwarzer Löcher an der Planck-Skala liefert das EABC-Modell exakte, unumstößliche Relationen. Das Framework postuliert die exakte Identität zwischen der mikroskopischen, konformen Entropie des starr modulierten Leech-Resonators ($S_{\text{Leech}}$) und der klassischen makroskopischen Bekenstein-Hawking-Entropie ($S_{\text{BH}}$) eines schwarzen Lochs:
\begin{equation}
    \mathbf{S_{\text{BH}} \equiv S_{\text{Leech}}}
\end{equation}

Setzt man die gravitative Flächenformel von Bekenstein-Hawking
\begin{equation}
    S_{\text{BH}} = \frac{A}{4\ell_P^2}
\end{equation}
mit der über die Cardy-Formel berechneten asymptotischen Entropie des $c=24$ Leech-Vakuums ($S_{\text{Leech}} \approx 4\pi\sqrt{E}$) gleich, resultiert die holografische Kopplungsbedingung:
\begin{equation}
    \mathbf{\frac{A}{4\ell_P^2} \equiv 4\pi \sqrt{E}}
\end{equation}

Diese Formel koppelt die makroskopische geometrische Horizontfläche $A$ der Einstein-Metrik bijektiv und verlustfrei an das mikroskopische Quantenspektrum $E$ des Resonators. Sie beweist, dass die Horizontfläche kein kontinuierliches Kontinuum darstellt, sondern das makroskopische Limit eines Gitters ist, dessen elementare Flächenquanten exakt durch die $196.560$ Minimalvektoren des Leech-Gitters moduliert werden.

\section{Frame Dragging und die Bändigung rotatorischer Weyl-Moden}

In der klassischen Allgemeinen Relativitätstheorie entsteht der Lense-Thirring-Effekt (Frame Dragging) durch Off-Diagonal-Terme $g_{0i}$ der Metrik, die auf den Drehimpuls $J$ zurückzuführen sind. Da der oktonionische Raum $\mathbb{O} \cong \mathbb{R}^8$ euklidisch ist, existiert dort kein direkter Lorentz-Zeit-Raum-Off-Diagonalterm.

Im EABC-Modell entsprechen rotatorische Anregungen den Generatoren außerhalb der $\operatorname{Spin}(7)$-Holonomie (den Weyl-Rauschmoden). Ein kontinuierliches Schleppen der Raumzeit wird jedoch durch die diskreten Symmetrien des Leech-Gitters und die Fano-Automorphismen topologisch quantisiert und unterbunden, wodurch rotatorische Fluktuationen in diskrete Sektoren gezwungen werden.

\section{Topologische Quantisierung der Perihel-Präzession an Gate $G_5$}

Die Kontrolle kontinuierlicher orbitaler Drifts lässt sich anhand des dreistufigen EABC-Modells für die Perihel-Präzession rigoros demonstrieren:

\subsection*{Stufe 1: Die geodätische Baseline}
Ausgangspunkt ist die Fermat-Primzahl im split-Sektor $p = 5 \equiv 1 \pmod 4$. Über die Gaußsche Zerlegung in $\mathbb{Z}[i]$ als $5 = (2+i)(2-i)$ ergibt sich die fundamentale Periodenlänge der geschlossenen Geodäte auf der Modulfläche $\operatorname{PSL}(2,\mathbb{Z})\backslash\mathbb{H}^2$:
\begin{equation}
    \ell(5) = \ln(5) \approx \mathbf{1{,}609437}
\end{equation}

\subsection*{Stufe 2: Spektraldynamik und Weyl-Dämpfung}
Ein unphysikalischer orbitaler Drift $A(\tau)$ unterliegt der spektralen Dämpfung im Co-Sektor mit der EABC-Grenzrate $\lambda_{\text{Weyl}} = 3/16 = 0{,}1875$. Nach einem vollen Umlauf $\tau = \ln(5)$ reduziert sich die Amplitude exponentiell:
\begin{equation}
    \frac{A(\ell(5))}{A(0)} = e^{-0{,}1875 \cdot \ln(5)} = 5^{-3/16} \approx \mathbf{0{,}7395}
\end{equation}

\subsection*{Stufe 3: Topologisches Phasen-Locking via Ptolemäus}
Um einen verbleibenden kontinuierlichen Rest-Drift vollständig einzufrieren, greift das ptolemäische Stiffness-Constraint auf dem Einheitskreis $U(1)$. Vier Bahnpunkte mit den Fermat-Gauss-Phasenwinkeln $\theta_k = k \cdot \arctan(1/2)$ spannen ein zyklisches Sehnenviereck auf. Der Satz von Ptolemäus fordert für die euklidischen Abstände:
\begin{equation}
    d_{13} \cdot d_{24} \equiv d_{12} \cdot d_{34} + d_{14} \cdot d_{23}
\end{equation}
Numerisch kollabiert diese Identität exakt zu:
\begin{equation}
    \mathbf{0{,}8000 \equiv 0{,}2111 + 0{,}5889 \equiv 0{,}8000}
\end{equation}

\subsection*{Physikalisches Fazit}
Das geometrische Korsett des Einheitskreises ist vollkommen unelastisch. Jede infinitesimale kontinuierliche Präzession $\delta\theta$ würde das Sehnenviereck verzerren und die ptolemäische Relation verletzen. Da der Gitter-Stabilisator diese algebraische Invarianz erzwingt, wird eine kontinuierliche Abweichung geometrisch blockiert. Die Perihel-Präzession ist somit topologisch auf diskrete Weyl-Sektoren quantisiert, wodurch die metrizierte Integrität des Vakuums gewahrt bleibt.

\end{document}